% !TeX root = main.tex
% Converted from chapter-one-revision-2026-10-17.md; source notes retained separately.
\chapter{Introduction}\label{chap:1}

\section{Background of the Study}\label{sec:1-1}

Student records provide information needed for departmental administration throughout a student's academic history. These records include identifying information, academic results and documents that support administrative decisions. Their value depends on whether authorised personnel can locate the correct information, interpret it in its proper context, and update it without losing its association with the student concerned. Filing is therefore more than the physical storage of documents; it involves organising records so that they remain retrievable, updatable and usable over time.

Gill~\cite[pp. 3-6]{gill1993} describes records management through the creation, use, storage and subsequent disposition of recorded information, with retrieval serving as an important test of a filing system's usefulness. These principles provide a foundation for departmental student filing. Records must be classified and indexed, linked to the appropriate student, and handled through defined procedures. Converting a paper document into an electronic file addresses its storage format, but effective digital filing also requires suitable metadata, retrieval arrangements and controls over access and correction.

In the Department of Electrical and Electronic Engineering, Obafemi Awolowo University, the project team has observed that student records are predominantly maintained in hard-copy form. Locating the required records can be difficult, and reviewing, updating and correcting information requires considerable effort. This observation motivates the investigation of a web-based filing system that can organise student information and supporting documents within an accessible departmental application.

A digital student filing system has two related forms of information to manage. Structured records contain fields such as student identifiers, biodata, academic periods and results that can be validated, queried and processed. Uploaded documents preserve supporting material, including scanned records, and require descriptive information that makes them identifiable and retrievable. These forms should be associated with the same student record while retaining their different storage, correction and access requirements. Uploading an image of a document does not by itself populate the structured database or make every word in the image searchable.

Web-based academic information systems provide an established approach to coordinating records and user responsibilities. Aliyu~\cite[pp. 2-3, 20-26]{aliyu2015} describes a university academic information system with integrated data and functions for different categories of users. Such work supports consideration of a departmental application in which the interface, application services and persistent records have defined responsibilities. A web interface can provide a common point of access, while server-side rules determine which information each user may view or change.

Records management also requires a distinction between locating a student's information and confirming the identity of a person presenting for a departmental activity. A name or matriculation number can identify a claimed record, but an identity-dependent workflow may require additional evidence that the person presenting is associated with that record. Fingerprint biometrics provides one means of obtaining such evidence through enrolment, acquisition and comparison against an enrolled reference. Recognition performance depends on the acquired sample, the reference data, the matching procedure and the operating conditions; it should therefore be evaluated rather than assumed.

Biometric links to student information are already represented in the literature. Omorogiuwa, Nwelih and Aweh~\cite{omorogiuwa2023} describe a fingerprint-based student information retrieval system, while Adetiba et al.~\cite{adetiba2013} describe online biometric class attendance management. These studies establish relevant prior approaches. The present project concerns their application within a departmental system that combines the management of structured records and uploaded documents with a defined fingerprint-assisted workflow.

The project therefore brings together a web-based student filing application and a fingerprint biometric subsystem. The filing application is responsible for organising and managing records. The biometric subsystem supplies identity evidence for selected activities, including the documented attendance workflow. Their integration must associate identity outcomes with the correct student record while preserving the permissions and rules governing that activity.

\section{Statement of the Problem}\label{sec:1-2}

Student records in the Department of Electrical and Electronic Engineering, Obafemi Awolowo University, are predominantly maintained as hard copies. As observed by the project team, these files are difficult to retrieve, and considerable effort is required to review, update and correct the information they contain. Access to the required information consequently depends on locating and handling the relevant physical records. This creates a practical records-management problem for departmental activities that require an accurate and up-to-date view of a student's information.

Addressing this problem requires more than storing electronic copies of existing papers. Structured student and academic information must be organised consistently, while uploaded documents must be indexed and associated with the correct student. Authorised users need to retrieve both forms of information and make controlled corrections without confusing students, academic periods or document versions. The first problem addressed by this project is therefore how to organise and manage these records within a coherent web-based departmental filing system.

A related problem arises in activities that require confirmation of a student's identity, such as attendance recording. The system must associate the person presenting with the appropriate student record before that identity outcome is used within the activity. Fingerprint recognition can support this association, but the recognition outcome must be integrated with the application's student identifiers, permissions and workflow rules. Capturing a fingerprint or receiving a matching result alone does not establish all the conditions required for a valid attendance record or unrestricted access to student information\footnote{}.

The project consequently addresses two connected requirements: the management of structured student records and supporting documents, and the integration of fingerprint-based identity recognition with selected departmental workflows. The engineering problem is to coordinate these functions so that records can be retrieved and updated appropriately and biometric outcomes can be linked to the correct student record under defined operating conditions.

\section{Research Questions}\label{sec:1-3}

\begin{enumerate}
\item How can a web-based departmental filing system be designed and implemented to organise structured student and academic records alongside uploaded documents, while supporting accurate retrieval, controlled updating and authorised access?
\item To what extent can a fingerprint biometric subsystem be integrated with that system to associate a presented fingerprint with the appropriate student record and support reliable, defined identity-dependent workflows?
\end{enumerate}

\section{Aim and Objectives}\label{sec:1-4}

The aim of this project is to design and implement an integrated fingerprint biometric and web-based student management system for organising departmental student records and supporting student identification within selected departmental workflows.

The specific objectives are to:

\begin{enumerate}
\item Analyse the existing departmental filing workflow and establish the functional and access requirements for student records, supporting documents and the selected biometric workflow.
\item Design a data model that associates student information, academic records, document metadata and biometric identity references through stable student identifiers and appropriate integrity rules.
\item Implement a web-based application for the controlled creation, updating and retrieval of structured records, together with the upload and retrieval of supporting documents.
\item Implement the fingerprint terminal functions required for acquisition, enrolment or recognition, user feedback and preservation of capture evidence under the selected architecture.
\item Integrate biometric identity outcomes with the relevant student records and selected attendance workflow, including the handling of unresolved identities, duplicate events and interrupted processing.
\item Evaluate the records functions, biometric subsystem and their integration using defined functional, recognition, timing and reliability criteria.
\end{enumerate}

\section{Significance of the Study}\label{sec:1-5}

The project is intended to provide a more organised approach to departmental student filing by associating structured information and supporting documents with the correct student record. This can support authorised personnel in locating relevant information and carrying out controlled updates and corrections. The intended benefit is directly related to the reported difficulty of working with predominantly hard-copy files.

The fingerprint subsystem adds a means of establishing an identity association within selected departmental activities. Its value depends on the correctness of that association and the rules applied afterwards. Evaluating the records and biometric functions together can show both the usefulness and practical limitations of the integrated prototype.

The study also provides engineering documentation of the records model, web application, fingerprint scanner terminal, communication arrangements and evaluation procedures. Its results can inform further development and departmental assessment. Claims of reduced retrieval time, reduced workload or improved recognition performance will be based on the measurements actually obtained rather than assumed from the adoption of web or biometric technology.

\section{Scope of the Project}\label{sec:1-6}

The project concerns a departmental prototype comprising a web-based student filing and management application and a fingerprint biometric subsystem. Its records scope includes structured student and academic information and uploaded supporting documents. The document functions require categorisation, descriptive metadata, association with the appropriate student, controlled upload and retrieval, and defined arrangements for corrections or replacement versions.

The document scope includes scanned records, photographs and other uploaded files that support the structured student information. It avoids treating the web application as a fully fledged student information system, all student results are still a product of the pre-existing university e-portal. The prototype does not attempt to make every word in an uploaded document searchable, nor does it attempt to convert all existing departmental records into electronic form. It is instead bounded by the selected student record set, the documents actually uploaded and the implemented filing functions.

The application includes access rules for the selected departmental users. Access to a record or an administrative action depends on the requesting user's permissions and relevant assignment or ownership rules. A fingerprint result supplies identity evidence for the selected workflow; it does not replace staff authentication or automatically grant access to academic information.

The biometric scope includes the terminal functions and identity-binding arrangements needed for the selected attendance application. Acquisition, matching, evidence transfer, identity resolution and attendance recording are distinct operations whose implementation and evaluation must be reported separately. Where deferred capture is used, an unresolved capture remains pending evidence until the required identity and workflow checks have been completed.

Evaluation covers the correctness of records and document operations, permitted and denied access, biometric decisions under the stated protocol, processing times and the integration cases actually implemented. The prototype is bounded by the selected hardware, software, record set and test conditions. Large-scale institutional deployment and biometric modalities other than fingerprint recognition are outside this prototype's scope.

\section{Definition of Key Terms}\label{sec:1-7}

% !TeX root = ../main.tex
\begingroup
\singlespacing%
\setlength{\tabcolsep}{4pt}
\begin{longtable}{@{}>{\RaggedRight\arraybackslash}p{\dimexpr0.280\linewidth-2\tabcolsep\relax}>{\RaggedRight\arraybackslash}p{\dimexpr0.720\linewidth-2\tabcolsep\relax}@{}}
\caption{Definition of key terms}\label{tab:1-terms}\\
\toprule
\textbf{Term} & \textbf{Meaning in this project} \\
\midrule
\endfirsthead%
\toprule
\textbf{Term} & \textbf{Meaning in this project} \\
\midrule
\endhead%
\midrule
\multicolumn{2}{r}{Continued on next page}\\
\endfoot%
\bottomrule
\endlastfoot%
Student filing system & The organisation and controlled management of structured student records and uploaded supporting documents. \\ \addlinespace
Structured record & Information held in defined database fields and relationships, such as student identifiers, academic periods and result entries. \\ \addlinespace
Supporting document & An uploaded digital file, including a scanned record, linked to a student and described by metadata used for filing and retrieval. \\ \addlinespace
Student registration & Creation or maintenance of a student's departmental record. It is distinct from course registration, account creation and biometric enrolment. \\ \addlinespace
Biometric enrolment & Acquisition and storage of a fingerprint reference and its association with the appropriate student identity. \\ \addlinespace
Fingerprint verification & Comparison of a new fingerprint sample against a reference for a claimed identity, commonly a 1:1 operation. \\ \addlinespace
Fingerprint identification & Search of a new fingerprint sample against an enrolled reference gallery to resolve an identity, commonly a 1:N operation. \\ \addlinespace
Authorisation & The application's decision about whether an authenticated user may perform a particular action on a record or other resource. \\ \addlinespace
Pending capture & Saved fingerprint evidence that has not yet completed the identity, timing and workflow checks required for an attendance decision. \\ \addlinespace
\end{longtable}
\endgroup


